\section{Methodology}

We test whether benchmark concentration creates epistemic lock-in through a six-hypothesis framework that traces the mechanism chain from concentration measurement through citation propagation to temporal dynamics. Our data comes from Papers With Code, covering 12,600 papers from NeurIPS, ICML, and ICLR spanning 2018--2024.

\subsection{Data Collection and Coverage}

Papers With Code provides structured information linking papers to the datasets they use for evaluation. We extract all entries for papers appearing at NeurIPS, ICML, or ICLR during our study period, yielding 12,600 papers with benchmark annotations. We estimate coverage exceeds 80\% of accepted papers at these venues based on comparing PWC entry counts against official venue acceptance statistics from conference proceedings.

For each paper, we record the set of datasets used for evaluation, venue, year, and citation links to other papers in the corpus. Citation relationships are operationalized through shared task annotations in Papers With Code, which serves as a proxy for intellectual influence: papers evaluating on the same tasks naturally cite each other. This task co-occurrence measure provides computational tractability while capturing the benchmark propagation mechanism of interest.

\subsection{Concentration Operationalization}

We measure benchmark concentration using the Herfindahl-Hirschman Index (HHI). For a given venue-year, let $N$ be the number of papers and let $n_d$ be the count of papers using dataset $d$. The share for dataset $d$ is $s_d = n_d / N$. HHI is then:
\begin{equation}
\text{HHI} = \sum_d s_d^2
\end{equation}

This metric equals the probability that two papers drawn uniformly at random from the venue-year evaluate on at least one common dataset. HHI ranges from $1/D$ (uniform distribution across $D$ datasets) to 1 (all papers use the same dataset). Our observed HHI values range from 0.007 to 0.046.

\subsection{Analytical Framework}

Our analysis traces the mechanism chain from concentration measurement to temporal dynamics. We first establish that HHI provides a valid concentration measure by verifying it can be computed across all venue-years and correlates with simpler metrics like top-5 dataset share. We then test whether concentration predicts individual paper behavior: do papers adopt standard benchmarks more readily when venue-level concentration is already high?

The propagation mechanism is tested through citation network analysis. If benchmark standards spread through intellectual influence, papers should share more evaluation overlap with the papers they cite than with random papers. Finally, we test for temporal lock-in directly: does concentration in year $t-1$ reduce diversity in year $t$? This is the critical test distinguishing co-evolution from lock-in.

\subsection{Panel Regression Specification}

The temporal lock-in test (H-M5) is our primary contribution. We estimate:
\begin{equation}
\text{Entropy}_{v,t} = \alpha + \beta \cdot \text{HHI}_{v,t-1} + \gamma_v + \delta_t + \epsilon_{v,t}
\end{equation}
where $\gamma_v$ represents venue fixed effects and $\delta_t$ represents year fixed effects. Significant negative $\beta$ would indicate lock-in; we additionally conduct Granger causality tests.
